\begin{solapartat}
\punts{0,2} $E = k\dfrac{q}{r^2}$

\medskip
\punts{0,4} $\dfrac{E_B}{E_A} = \dfrac{k\dfrac{q_U}{r_B^2}}{k\dfrac{q_U}{r_A^2}} = \dfrac{r_A^2}{r_B^2}$

\medskip
\punts{0,4} $\dfrac{E_B}{E_A} = \dfrac{\left(0,008\times 10^{-9}\right)^2}{\left(0,008\times 10^{-12}\right)^2} = 10^{6}$

\medskip
És a dir, el mòdul del camp elèctric en B es un milió de vegades més gran que el mòdul del camp elèctric en A.
\end{solapartat}

\begin{solapartat}
\punts{0,2} Cal que es conservi l'energia mecànica: $K_A + U_A = K_B + U_B$. Com a mínim $K_B = 0 \rightarrow K_A + U_A \geq U_B \rightarrow K_A \geq U_B - U_A$.

Podria negligir-se la $U_A$ i obtenir legítimament el mateix resultat, però caldria incloure alguna justificació explícita (per exemple ``tenint present que $|U_A| \ll |U_B|$\ldots'')

\medskip
\punts{0,4} $K_A \geq k\dfrac{q_\alpha q_U}{r_B} - k\dfrac{q_\alpha q_U}{r_A} = kq_\alpha q_U\left(\dfrac{1}{r_B} - \dfrac{1}{r_A}\right)$

\medskip
\punts{0,4}
$\begin{aligned}[t]
K_A &\geq 8,99\cdot 10^{9}\cdot\underbrace{2\cdot 1,6\times 10^{-19}}_{q_\alpha}\cdot\underbrace{92\cdot 1,6\times 10^{-19}}_{q_U}\cdot\left(\frac{1}{0,008\times 10^{-12}} - \frac{1}{0,008\times 10^{-9}}\right)\\
K_A &\geq 5,29\times 10^{-12}\ \mathrm{J}
\end{aligned}$
\end{solapartat}
